% Proved after four complete actual-file audits. Global closure of the same bounded-period hereditary crown reduction.
\subsection{Global symplectic and orthogonal crowns over arbitrary quotients}

\begin{corollary}[Proved: global symplectic and orthogonal crown bounds]
\label{D:global-arbitrary-quotient-symplectic-orthogonal-crowns}
Let $X$ be a fixed normal proper connected surface over an
algebraically closed field of characteristic $p>0$, let
$T\subset X$ be finite, and put $\Gamma=\pi_1(X-T)$.
Let $S$ be any of the nonabelian simple groups, in precisely
the ranges and actual coefficient-field models of
Corollary~\ref{D:arbitrary-quotient-symplectic-orthogonal-crowns}.
Write $d,Q$ for its model dimension and field cardinality.
Let $L$ be a finite monolithic primitive group with socle
$N=S^t$, put $P=L/N$, and prescribe a continuous epimorphism
$\beta:\Gamma\to P$.

Let $A^{\rm sl}_{X,T,2}$ be the constant of
Theorem~\ref{D:global-finite-puncture-hereditary-semilinear-local-period}
with period bound two. Every crown quotient
$\Gamma\twoheadrightarrow L_k$ over $\beta$ satisfies
\[
 \begin{aligned}
 k&\le |N|^{c_{X,T}},\\
 d(L_k)&\le\max\{2,d(P),c_{X,T}+2\},\\
 c_{X,T}&=\max\{1,\lceil96A^{\rm sl}_{X,T,2}+21\rceil\}.
 \end{aligned}                                               \tag{D.GO.1}
\]
There is no tameness, resolution-extension or bounded
field-image assumption. The surface and finite omitted
set remain fixed; the constant is independent of the
classical rank, defining field, $t$ and prescribed
permutation quotient. The term $d(P)$ is retained.
\end{corollary}

\begin{proof}
Use the actual recognition proved in
Corollary~\ref{D:arbitrary-quotient-symplectic-orthogonal-crowns}.
Its normal inner-diagonal-field subgroup
$\mathcal A_0\subset\operatorname{Aut}(S)$ has index
at most six and an actual faithful embedding
\[
 \begin{gathered}
 \mathcal A_0\lhook\joinrel\longrightarrow
 \operatorname{PGL}_d(\mathbb F_Q)\rtimes
                  \operatorname{Gal}(\mathbb F_Q/\mathbb F_p),\\
 S\subset\operatorname{PGL}_d(\mathbb F_Q),
 \qquad d^2\log Q\le16\log|S|.
 \end{gathered}                                             \tag{D.GO.2}
\]
Its full inverse image consists of semisimilitudes
of one fixed nonzero prime-field bilinear or quadratic
form, with scalar multiplier $u^2$.
In particular the minus orthogonal row keeps
$Q=q^2$ and its actual field subgroup, exactly as in
that recognition proof. No additional family or
smaller-field model is introduced here.

Choose a simple socle factor. Let $H\subset\Gamma$
be the inverse image under $\beta$ of its point
stabilizer, so $[\Gamma:H]=t$. Let
$\mathcal B\subset\operatorname{Aut}(S)$ be its
stabilizer image in $L$, and let
$\eta:H\to\mathcal B/S$ be the fixed outer map.
Set
\[
 \mathcal B_0=\mathcal B\cap\mathcal A_0,
 \qquad H_0=\eta^{-1}(\mathcal B_0/S),
 \qquad a=[H:H_0]\le6 .                                  \tag{D.GO.3}
\]
Since $S$ is perfect, the field projection on
$\mathcal B_0$ factors through $\mathcal B_0/S$.
Thus $\eta$ fixes a single field character $\chi_0$
on $H_0$ for every stabilizer lift under consideration.
The connected cover corresponding to $H_0$ has degree
$at$ over $X-T$.

We verify the precise hypothesis of the global
hereditary theorem on this cover. For a stabilizer
lift $\theta:H\to\mathcal B$, choose set-theoretic
semilinear matrix lifts $\widetilde\theta(g)$ of
its restriction to $H_0$, and write
\[
 \widetilde\theta(g)\widetilde\theta(h)
       =c(g,h)\widetilde\theta(gh).
\]
The form multiplier $\lambda$ has the exact crossed
multiplicativity and scalar-square properties from
(D.BS.6). Consequently
\[
 c(g,h)^2=
 \frac{\lambda(\widetilde\theta(g))\,
       \chi_0(g)(\lambda(\widetilde\theta(h)))}
      {\lambda(\widetilde\theta(gh))} .                   \tag{D.GO.4}
\]
The full obstruction class in
$H^2(H_0,\mathbb F_Q^\times(\chi_0))$ is therefore
killed by two. This conclusion is obtained before
passing to any geometric Brauer group.

Use the fixed complete local sources of the global
hereditary theorem. For each connected factor of
the pullback of such a source to the degree-$at$
cover, restrict this class to its fundamental group,
and then to the kernel of the actual restricted
field character. Functorial restriction preserves
(D.GO.4). On that kernel the scalar band is constant,
so its ordinary coefficient class and its geometric
Brauer image are killed by two.
This is exactly the local period condition of
Theorem~\ref{D:global-finite-puncture-hereditary-semilinear-local-period}.
It does not require the restricted field image to
have the same order as the global image, any
injectivity of a local fundamental-group map,
or any comparison of Brauer groups in the opposite
direction. No ordinary Brauer class is assigned to
the nonconstant band on the degree-$at$ cover itself.

That theorem gives at most
$Q^{A^{\rm sl}_{X,T,2}atd^2}$ ambient projective
conjugacy classes for the restricted maps on $H_0$.
The exact nonsplit transport of
Lemma~\ref{D:nonsplit-automorphism-wreath-lifts}
and the orbit and extension argument of
Corollary~\ref{D:bounded-period-arbitrary-quotient-projective-semilinear-crowns}
are purely profinite and apply to the present
$\Gamma$. We spell out the resulting numerical
bound. An ambient conjugacy orbit contains at most
$|\operatorname{PGL}_d(\mathbb F_Q)|\le Q^{d^2}$
actual homomorphisms, hence at most that many
$S$-classes of $\mathcal B_0$-valued restrictions.
Using (D.GO.2), their number is at most
\[
 Q^{(A^{\rm sl}_{X,T,2}at+1)d^2}
       \le |S|^{16(A^{\rm sl}_{X,T,2}at+1)}.             \tag{D.GO.5}
\]
Fix an actual representative of each restricted
$S$-class. Conjugating by $S$ normalizes every full
lift in that class to this restriction and preserves
its prescribed outer map. The images of $a-1$
nonidentity coset representatives of $H_0$ in $H$
then have at most $|S|^{a-1}$ choices, because each
lies in one prescribed $S$-coset. The restriction
and these images determine the whole homomorphism.
This argument uses actual $S$-normalization, not
an action of the whole stabilizer on the ambient
projective space.

Surjectivity onto the crown power forces its
coordinate lifts to represent distinct stabilizer
$S$-classes, by the same diagonal-graph argument
in the exact transport lemma. Hence
\[
 \begin{aligned}
 k&\le |S|^{16(A^{\rm sl}_{X,T,2}at+1)+a-1}\\
  &\le |S|^{t(96A^{\rm sl}_{X,T,2}+21)}
   \le |N|^{c_{X,T}}.
 \end{aligned}                                              \tag{D.GO.6}
\]
The exact nonabelian crown inequality used in the
preceding crown corollary, from the proof of
\citet[Proposition~2.6(ii), equations~(2.1)--(2.2),
pp.~11--13]{LucchiniMarionTracey2020D}, then gives
$d(L_k)\le\max\{2,d(P),c_{X,T}+2\}$.
This establishes (D.GO.1) without bounding $d(P)$.
\end{proof}

\begin{remark}[Scope of the global recognition closure]
This applies only to a fixed normal proper surface
with a finite omitted set. No hereditary theorem
for a divisorial boundary is used or asserted.
The named simple families and their low-rank
exclusions are exactly those of the local
recognition corollary. In particular the unitary
Hermitian scalar exponent $q+1$ and unbounded-centre
type $A$ are not included. This is a global closure
of the same bounded-period hereditary crown
reduction, not a separate movement or a finite-
generation theorem for $\Gamma$.
\end{remark}
